\section{问题二：商品关联分析与类别聚合预测}

问题二关注同一门店中不同零食销量之间的联系，并要求整合同类零食预测未来 7 天总销量。设 $y_{t,p}$ 表示第 $t$ 天商品 $p$ 的正向销量，$c(p)$ 表示商品所属类别。

\subsection{商品销量矩阵与相关性}

对每个门店构造日期 $\times$ 商品销量矩阵：
\begin{equation}
Y_s=(y_{t,p}),\quad t=1,\ldots,n,\ p=1,\ldots,m.
\end{equation}
商品 $p$ 与商品 $q$ 的 Pearson 相关系数为
\begin{equation}
r_{pq}=
\frac{\sum_t(y_{t,p}-\bar y_p)(y_{t,q}-\bar y_q)}
{\sqrt{\sum_t(y_{t,p}-\bar y_p)^2}\sqrt{\sum_t(y_{t,q}-\bar y_q)^2}}.
\end{equation}
为避免单一 Pearson 相关造成过度解释，本文补充 Spearman 相关、去星期效应 Pearson 和门店内相关稳定性检查。商品关联结论只解释为同步波动，不写成因果或替代关系。

\begin{figure}[H]
\centering
\includegraphics[width=0.95\textwidth]{q2_overall_product_correlation_heatmap.png}
\caption{商品日销量相关性热力图}
\label{fig:q2_corr}
\end{figure}

全部门店汇总层面共计算 $12$ 个商品代码之间的 66 对 Pearson 相关系数。本文将 $r\ge 0.50$ 记为强正相关，将 $|r|\le 0.10$ 记为弱线性相关。结果中有 2 对强正相关商品对，见表 \ref{tab:q2_strong_pairs}；有 25 对弱相关商品对，见表 \ref{tab:q2_weak_pairs}。

\begin{table}[H]
\centering
\caption{强正相关商品对}
\label{tab:q2_strong_pairs}
\begin{tabular}{llrr}
\toprule
商品 A & 商品 B & 相关系数 & 同日均有销量天数 \\
\midrule
三养辣火鸡味拌面 & 乡吧哥蜜汁鸡翅 & 0.521 & 516 \\
三养辣火鸡味拌面 & 李子柒螺蛳粉335g & 0.506 & 400 \\
\bottomrule
\end{tabular}
\end{table}

\begingroup
\small
\setlength{\tabcolsep}{4pt}
\begin{longtable}{p{4.4cm}p{4.6cm}p{1.7cm}p{2.0cm}}
\caption{弱相关商品对（$|r|\le 0.10$）}
\label{tab:q2_weak_pairs}\\
\toprule
商品 A & 商品 B & 相关系数 & 同日均有销量天数 \\
\midrule
\endfirsthead
\toprule
商品 A & 商品 B & 相关系数 & 同日均有销量天数 \\
\midrule
\endhead
老板仔香脆海苔卷（原味） & 洽洽早餐每日坚果（燕麦片） & -0.063 & 11 \\
大丰古法瓜子 & 洽洽早餐每日坚果（燕麦片） & -0.055 & 16 \\
韩国海牌海苔 / 韩国海飘海苔 & 洽洽早餐每日坚果（燕麦片） & -0.026 & 17 \\
特仑苏纯牛奶 & 李子柒螺蛳粉335g & -0.025 & 254 \\
特仑苏纯牛奶 & 红牛 & -0.018 & 390 \\
特仑苏纯牛奶 & 洽洽早餐每日坚果（燕麦片） & -0.016 & 9 \\
百事可乐330ml & 洽洽早餐每日坚果（燕麦片） & -0.009 & 17 \\
李子柒螺蛳粉335g & 洽洽早餐每日坚果（燕麦片） & -0.004 & 16 \\
洽洽早餐每日坚果（燕麦片） & 红牛 & 0.018 & 17 \\
三养辣火鸡味拌面 & 洽洽早餐每日坚果（燕麦片） & 0.019 & 14 \\
百事可乐330ml & 特仑苏纯牛奶 & 0.020 & 375 \\
乡吧哥蜜汁鸡翅 & 特仑苏纯牛奶 & 0.020 & 380 \\
特仑苏纯牛奶 & 大丰古法瓜子 & 0.020 & 377 \\
三养辣火鸡味拌面 & 特仑苏纯牛奶 & 0.023 & 288 \\
乡吧哥蜜汁鸡翅 & 洽洽早餐每日坚果（燕麦片） & 0.024 & 17 \\
韩国海牌海苔 / 韩国海飘海苔 & 特仑苏纯牛奶 & 0.028 & 374 \\
洽洽早餐每日坚果（燕麦片） & 旺仔牛奶 & 0.034 & 17 \\
特仑苏纯牛奶 & 百事可乐600ml & 0.036 & 391 \\
韩国海牌海苔 / 韩国海飘海苔 & 红牛 & 0.037 & 699 \\
洽洽早餐每日坚果（燕麦片） & 百事可乐600ml & 0.041 & 17 \\
特仑苏纯牛奶 & 旺仔牛奶 & 0.045 & 368 \\
李子柒螺蛳粉335g & 红牛 & 0.054 & 491 \\
大丰古法瓜子 & 红牛 & 0.065 & 699 \\
百事可乐330ml & 大丰古法瓜子 & 0.085 & 679 \\
老板仔香脆海苔卷（原味） & 特仑苏纯牛奶 & 0.091 & 301 \\
\bottomrule
\end{longtable}
\endgroup

表 \ref{tab:q2_strong_pairs} 中两对强正相关都包含三养辣火鸡味拌面，更合理的解释是它们在历史日销量上同步波动，可能共同受到客流、即时餐食消费场景或促销陈列节奏影响。表 \ref{tab:q2_weak_pairs} 中，多数弱相关对涉及低销量商品或消费场景差异较大的商品；这说明在当前附件数据下，不能把相关性弱或略为负相关直接解释为替代关系。

\subsection{同类零食整合}

类别聚合按附件 \texttt{category} 字段进行：
\begin{equation}
Y_{t,g}=\sum_{p:c(p)=g}y_{t,p}.
\end{equation}
该字段具有业务含义，因此作为“同类零食”的正式聚合依据。相关性聚类和销量规模分组可作为补充探索，但不替代附件业务类别。

\begin{figure}[H]
\centering
\includegraphics[width=0.94\textwidth]{q2_category_daily_trend.png}
\caption{各类别日销量历史趋势}
\label{fig:q2_category_trend}
\end{figure}

\subsection{类别预测与误差比较}

本文比较“单品预测后按类别加总”和“类别聚合后直接预测”。验证结果见表 \ref{tab:q2_metrics}，口径为 2022 年 3 月日滚动一步验证。

\begin{table}[H]
\centering
\caption{问题二类别预测策略误差比较（日滚动一步验证）}
\label{tab:q2_metrics}
\begin{tabular}{llrrr}
\toprule
方法 & 模型 & MAE & RMSE & WAPE \\
\midrule
类别聚合后直接预测 & 简单指数平滑 & 7.853 & 13.933 & 34.74\% \\
单品预测后按类别加总 & 简单指数平滑 & 7.873 & 13.949 & 34.83\% \\
类别聚合后直接预测 & 移动平均(7日) & 8.246 & 14.174 & 36.48\% \\
类别聚合后直接预测 & 同星期均值(近8周) & 8.851 & 15.057 & 39.15\% \\
\bottomrule
\end{tabular}
\end{table}

因此，问题二主方案采用“附件类别聚合后直接预测 + 简单指数平滑”。方法复核中，相关性聚类在同一日滚动一步验证下 WAPE=31.02\%，销量规模分组 WAPE=24.24\%，但前者业务类别含义弱，后者预测对象变成高/中/低销量组，均不替代附件类别主方案。

\subsection{问题二结果}

商品关联分析表明，全部门店汇总层面存在 2 对强正相关商品对和 25 对弱相关商品对。强正相关只能说明历史日销量同步波动，可能共同受客流、消费场景或促销节奏影响，不能证明商品之间存在直接带动关系。负相关候选的相关系数绝对值较小，当前数据未发现强替代关系证据。

类别预测结果显示，未来 7 天预测最高的类别为包装散称，预测 389.613；其次为功能饮料，预测 286.719；碳酸饮料预测 247.326。类别聚合可以降低单品零销量和偶然大单带来的波动，但会损失同一类别内部商品差异。
